% ====================================================================
% KONFIGURASI HEADER, FOOTER, DAN PENOMORAN HALAMAN
% ====================================================================
% File ini mengatur penomoran halaman dan header/footer dokumen.
%
% Ketentuan Penomoran Juknis UNEJ:
% - Bagian Awal (Front Matter): Roman kecil (i, ii, iii, ...) di tengah bawah
% - Bagian Inti (Main Body): Angka Arab (1, 2, 3, ...) di kanan atas
% - Halaman pertama setiap BAB: Angka di tengah bawah (tanpa header)
% ====================================================================

% === STYLE UNTUK HALAMAN REGULER ===
% Nomor halaman di kanan atas, tanpa garis header
\fancypagestyle{plain}{%
	\fancyhf{}%
	\fancyfoot[C]{\thepage}%
	\renewcommand{\headrulewidth}{0pt}%
	\renewcommand{\footrulewidth}{0pt}%
}

% Style untuk halaman biasa bagian utama: nomor di kanan atas
\fancypagestyle{fancy}{%
	\fancyhf{}%
	\fancyhead[R]{\thepage}%
	\renewcommand{\headrulewidth}{0pt}%
	\renewcommand{\footrulewidth}{0pt}%
}





% ====================================================================
% CATATAN IMPLEMENTASI DI main.tex:
% 1. Bagian Front Matter harus menggunakan:
%    \pagenumbering{roman}
%    \pagestyle{plain}
%    Untuk penomoran Romawi (i, ii, iii) di tengah bawah
%
% 2. Bagian Main Body harus menggunakan:
%    \pagenumbering{arabic}
%    \setcounter{page}{1}
%    \pagestyle{fancy}
%    Untuk penomoran Arab (1, 2, 3) di kanan atas
% ====================================================================
